\section{Carbon Exposure with Unmeasured Power}
\label{sec:accounting}

Fix one machine--workload pair; its indices are suppressed below. The mean
complete-node power assigned to scale $n$ is
\begin{equation}
 \widehat P_n=P_0\kappa s_n,\qquad
 s_n=\rho+(1-\rho)\eta_n,\quad
 \eta_n=\frac{W_4}{W_n},
 \label{eq:power}
\end{equation}
where $W_n=nU_{\rm profile}/q_n$, $\kappa>0$, $s_4=1$, and
$0\leq\rho\leq1$. The reference $P_0=1$\,kW fixes a normalization unit; $\kappa$ absorbs the
unknown workload-specific node-power level. Neither is a reported A800 server
measurement. We report within-panel carbon ratios and node-hours, rather than
absolute carbon or specification-based node-power estimates.
The workload coefficient $\kappa$ sets the four-node power level. The scale
factor interpolates between constant per-node power ($\rho=1$) and a
constant-useful-work dynamic-energy limit ($\rho=0$, without overhead).
$\eta_n$ is a scaling-efficiency ratio, not measured utilization, and $\rho$
is not the CPU fraction of node power. We use this family as a sensitivity
model; its parameter interval is a declared scenario, not a confidence interval.
At constant CI and without overhead, Equation~\ref{eq:power} gives total energy
$P_0\kappa[\rho W_\pi+(1-\rho)W_4^\star]$, where $W_\pi$ is the
execution's node-hours and $W_4^\star$ is the cost of doing all target updates
at four nodes. For $\rho>0$, changing this parameter cannot reverse an energy
ranking by node-hours. It changes the magnitude of savings in this limit,
not which execution uses less energy.

For an execution log $\pi$, let $\mathcal I_{\pi,n}$ contain all intervals
allocated to the target at scale $n$, including restart and checkpoint. With
realized average grid intensity $c(t)$, define the scale-resolved exposure
\begin{equation}
 L_{\pi,n}=n\int_{\mathcal I_{\pi,n}}c(t)\,dt,
 \qquad
 C_\pi(\rho)=P_0\kappa\sum_n s_n L_{\pi,n}.
 \label{eq:exposure}
\end{equation}
Power is in kW and time in hours. For the primary EIA scenario, intensity is
in kgCO$_2$/kWh; absolute $C_\pi$ remains unknown without $\kappa$.
Non-CO$_2$ greenhouse gases are
excluded. Replacing $c(t)$ with 1 yields IT energy; summing
allocated node time yields node-hours without a power assumption. Applying
the same mean $\widehat P_n$ to all allocated phases is an explicit
approximation. Phase durations are retained in the log so its effect can be
examined without claiming phase-specific power measurements.

Writing $A_\pi=\sum_nL_{\pi,n}$ and
$B_\pi=\sum_n\eta_nL_{\pi,n}$ gives
\begin{equation}
 C_\pi(\rho)=P_0\kappa
              [B_\pi+\rho(A_\pi-B_\pi)].
 \label{eq:affine}
\end{equation}
\begin{proposition}[Endpoint comparison]
\label{prop:endpoints}
For fixed policies and execution logs, $C_\pi(\rho)\leq C_b(\rho)$ for every
$\rho\in[\rho_-,\rho_+]$ if and only if this inequality holds at both
endpoints. The same statement holds for expected carbon under a fixed
distribution of arrivals and policy randomness.
\end{proposition}
\noindent\emph{Proof.}
The difference is affine in $\rho$ by Equation~\ref{eq:affine}; its maximum
on a closed interval occurs at an endpoint. Expectation preserves the affine
form.\hfill$\square$

The common positive factor $P_0\kappa$ cancels from within-panel
ratios and rankings. No workload coefficient sweep is needed for these
comparisons, although unknown coefficients preclude pooling absolute carbon
across workloads. For a positive affine reference $C_{\rm ref}(\rho)$, the
ratio $\mathbb E[C_\pi(\rho)]/C_{\rm ref}(\rho)$ is linear-fractional and
has a derivative of constant sign (or zero). Its worst value also occurs at
an endpoint. Policies must remain fixed while computing this result: retraining
at each $\rho$ would evaluate different policies, not certify one policy.

For powers outside this family, Equation~\ref{eq:exposure} remains useful.
If each scale exposure is no greater than a baseline's, dominance holds for
all positive scale-power coefficients. Otherwise a reversal is possible,
and the result must retain a power-model qualification. A more informative
local test allows $s_n=s_n^0(1+\delta_n)$ for $n\ne4$, with
$|\delta_n|\leq d<1$ and $s_4=1$. Writing
$\Delta L_n=L_{\pi,n}-L_{b,n}$, the worst normalized difference is exactly
\begin{equation}
 \max_{\delta}\frac{C_\pi-C_b}{P_0\kappa}
 =\sum_n s_n^0\Delta L_n
       +d\sum_{n\ne4}s_n^0|\Delta L_n|.
 \label{eq:box}
\end{equation}
This tests independent scale-power errors using the same logs, without
adding a power-model layer. Queueing draws no
allocated energy at this accounting boundary. We exclude embodied emissions,
idle and background-job power, and shared facility overhead. The estimated
quantity is job-attributed operational carbon, not avoided grid emissions.
